\chapter{Phân tích yêu cầu}

\section{Tổng quan hệ thống}
\subsection{Bối cảnh và mục tiêu}
\subsection{Phạm vi hệ thống}
\subsection{Đối tượng sử dụng}
\subsection{Các nhóm nghiệp vụ chính}

\section{Phân tích yêu cầu dữ liệu}
\subsection{Nhóm dữ liệu người dùng}
\subsection{Nhóm dữ liệu nhà cung cấp và tài nguyên}
\subsection{Nhóm dữ liệu dịch vụ}
\subsection{Nhóm dữ liệu đặt chỗ và sử dụng dịch vụ}
\subsection{Nhóm dữ liệu cộng đồng}
\subsection{Nhóm dữ liệu đánh giá và hỗ trợ}

\section{Phân tích các nghiệp vụ chính}
\subsection{Quản lý người dùng}
\subsection{Quản lý nhà cung cấp và tài nguyên}
\subsection{Quản lý dịch vụ}
\subsection{Đặt và sử dụng dịch vụ}
\subsection{Hoạt động cộng đồng}
\subsection{Đánh giá và xử lý khiếu nại}

\section{Đặc tả yêu cầu dữ liệu}

Phần này xác định các nhóm thông tin mà hệ thống cần lưu trữ và quản lý để hỗ trợ các nghiệp vụ chính của nền tảng. Việc đặc tả ở giai đoạn này tập trung vào ý nghĩa nghiệp vụ và loại thông tin cần quản lý, chưa xác định cấu trúc bảng, khóa hay kiểu dữ liệu.

\subsection{User}

Hệ thống cần quản lý thông tin người dùng tham gia và sử dụng nền tảng.

Các yêu cầu dữ liệu chính:

\begin{itemize}
\item Thông tin nhận diện và liên hệ của người dùng.
\item Thông tin tài khoản và trạng thái hoạt động.
\item Vai trò của người dùng trong hệ thống.
\item Thông tin phục vụ việc sử dụng dịch vụ, đặt chỗ và tham gia cộng đồng.
\item Lịch sử hoạt động cần thiết của người dùng như đặt dịch vụ, đánh giá, tham gia Workshop hoặc chia sẻ nội dung.
\end{itemize}

Hệ thống cần hỗ trợ các vai trò Photographer, Service Provider, Photography Expert và Administrator.

\subsection{Service Provider}

Hệ thống cần quản lý thông tin về các cá nhân hoặc tổ chức cung cấp dịch vụ phòng tối, phòng chụp, thiết bị và các dịch vụ sáng tạo.

Các yêu cầu dữ liệu chính:

\begin{itemize}
\item Thông tin của nhà cung cấp.
\item Thông tin liên hệ và địa điểm hoạt động.
\item Thông tin mô tả về đơn vị cung cấp dịch vụ.
\item Trạng thái hoạt động hoặc phê duyệt của nhà cung cấp.
\item Thông tin phục vụ quản lý các Creative Space, Resource và Service Package thuộc nhà cung cấp.
\item Thông tin cần thiết cho việc theo dõi hoạt động kinh doanh.
\end{itemize}

\subsection{Creative Space}

Hệ thống cần quản lý thông tin về các không gian sáng tạo được cung cấp trên nền tảng, chẳng hạn như darkroom và photography studio.

Các yêu cầu dữ liệu chính:

\begin{itemize}
\item Thông tin nhận diện và mô tả không gian.
\item Loại không gian.
\item Vị trí và thông tin liên quan đến địa điểm.
\item Diện tích và sức chứa.
\item Phong cách nghệ thuật.
\item Điều kiện ánh sáng.
\item Điều kiện thông gió và đặc điểm âm học.
\item Thời gian hoạt động.
\item Quy định và điều kiện sử dụng.
\item Thông tin về giá sử dụng.
\item Trạng thái và khả năng cung cấp dịch vụ.
\item Thông tin về các tiện ích và thiết bị được hỗ trợ.
\end{itemize}

Dữ liệu Creative Space phải hỗ trợ việc tìm kiếm, so sánh, kiểm tra khả dụng và đặt chỗ.

\subsection{Resource}

Hệ thống cần quản lý các tài nguyên được sử dụng hoặc cho thuê trong quá trình cung cấp dịch vụ.

Các yêu cầu dữ liệu chính:

\begin{itemize}
\item Thông tin nhận diện và mô tả tài nguyên.
\item Loại tài nguyên.
\item Số lượng tài nguyên.
\item Tình trạng hoạt động.
\item Trạng thái khả dụng.
\item Giá thuê hoặc chi phí sử dụng.
\item Thông tin về khả năng tương thích với các thiết bị hoặc dịch vụ khác.
\item Thông tin phục vụ theo dõi lịch sử sử dụng.
\item Thông tin phục vụ quản lý vòng đời tài nguyên.
\end{itemize}

Resource có thể bao gồm camera, lens, enlarger, film scanner, lighting system, tripod, background, thiết bị darkroom, hóa chất và photographic paper.

\subsection{Maintenance}

Hệ thống cần quản lý thông tin về hoạt động bảo trì đối với các Resource.

Các yêu cầu dữ liệu chính:

\begin{itemize}
\item Thông tin về tài nguyên cần bảo trì.
\item Loại và nội dung công việc bảo trì.
\item Thời gian dự kiến thực hiện.
\item Thời gian hoàn thành.
\item Chi phí bảo trì.
\item Trạng thái bảo trì.
\item Thông tin phục vụ theo dõi lịch sử bảo trì và vòng đời của tài nguyên.
\end{itemize}

Dữ liệu Maintenance phải hỗ trợ việc xác định tài nguyên đang được bảo trì để tránh phân bổ tài nguyên không khả dụng cho Reservation.

\subsection{Service Package}

Hệ thống cần quản lý các gói dịch vụ do Service Provider xây dựng và cung cấp.

Các yêu cầu dữ liệu chính:

\begin{itemize}
\item Thông tin tên và mô tả gói dịch vụ.
\item Nội dung và phạm vi của gói.
\item Giá và thời lượng sử dụng.
\item Trạng thái cung cấp.
\item Các Creative Space được bao gồm.
\item Các Resource được bao gồm.
\item Các dịch vụ hỗ trợ khác nếu có.
\end{itemize}

Service Package phải hỗ trợ các nhu cầu khác nhau như portrait photography, product photography, film developing, darkroom printing và educational workshops.

\subsection{Promotion}

Hệ thống cần quản lý các chương trình khuyến mãi được Service Provider cung cấp.

Các yêu cầu dữ liệu chính:

\begin{itemize}
\item Thông tin chương trình khuyến mãi.
\item Nội dung và điều kiện áp dụng.
\item Hình thức giảm giá.
\item Giá trị giảm giá.
\item Thời gian bắt đầu và kết thúc.
\item Trạng thái chương trình.
\end{itemize}

Dữ liệu Promotion được sử dụng để hỗ trợ việc áp dụng chính sách giá và khuyến mãi cho các dịch vụ được cung cấp trên nền tảng.

\subsection{Reservation}

Hệ thống cần quản lý thông tin về các yêu cầu đặt Creative Space, Resource hoặc Service Package của Photographer.

Các yêu cầu dữ liệu chính:

\begin{itemize}
\item Người thực hiện đặt dịch vụ.
\item Đối tượng và tài nguyên được đặt.
\item Khoảng thời gian đặt.
\item Trạng thái của Reservation.
\item Tổng giá trị đặt dịch vụ.
\item Thời điểm tạo yêu cầu đặt.
\item Thông tin cần thiết để kiểm tra tình trạng khả dụng.
\item Thông tin phục vụ theo dõi lịch sử đặt dịch vụ.
\end{itemize}

Hệ thống phải sử dụng dữ liệu Reservation để kiểm tra khả dụng và phát hiện xung đột lịch trước khi xác nhận đặt chỗ.

Sau khi Reservation được xác nhận, các tài nguyên liên quan phải được giữ cho khoảng thời gian đã đặt.

\subsection{Payment}

Hệ thống cần quản lý thông tin về các giao dịch thanh toán phát sinh từ hoạt động đặt dịch vụ.

Các yêu cầu dữ liệu chính:

\begin{itemize}
\item Reservation liên quan đến giao dịch.
\item Số tiền thanh toán.
\item Phương thức thanh toán.
\item Trạng thái giao dịch.
\item Thông tin nhận diện giao dịch.
\item Thời điểm thanh toán.
\item Thông tin phục vụ đối soát và theo dõi lịch sử thanh toán.
\end{itemize}

Dữ liệu Payment hỗ trợ việc quản lý giao dịch, theo dõi thanh toán và xử lý các vấn đề liên quan đến giao dịch.

\subsection{Service Session}

Hệ thống cần quản lý thông tin về quá trình sử dụng dịch vụ thực tế sau khi Reservation được xác nhận.

Các yêu cầu dữ liệu chính:

\begin{itemize}
\item Reservation tương ứng.
\item Thông tin check-in.
\item Thông tin check-out.
\item Thời lượng sử dụng thực tế.
\item Trạng thái phiên sử dụng.
\item Các Resource thực tế được sử dụng.
\item Các dịch vụ bổ sung đã sử dụng nếu có.
\end{itemize}

Service Session được sử dụng để phục vụ billing, đánh giá dịch vụ, theo dõi thiết bị và giải quyết tranh chấp.

\subsection{Review}

Hệ thống cần quản lý phản hồi của người dùng đối với dịch vụ hoặc tài nguyên đã sử dụng.

Các yêu cầu dữ liệu chính:

\begin{itemize}
\item Người thực hiện đánh giá.
\item Đối tượng hoặc dịch vụ được đánh giá.
\item Điểm đánh giá.
\item Nội dung nhận xét.
\item Thời điểm đánh giá.
\item Thông tin cần thiết để liên kết đánh giá với quá trình sử dụng dịch vụ.
\end{itemize}

Review được sử dụng để cung cấp thông tin phản hồi cho cộng đồng và hỗ trợ đánh giá chất lượng của Service Provider và dịch vụ.

\subsection{Community Content}

Hệ thống cần quản lý các nội dung được người dùng chia sẻ trong cộng đồng nhiếp ảnh phim.

Các yêu cầu dữ liệu chính:

\begin{itemize}
\item Thông tin về nội dung được chia sẻ.
\item Tiêu đề và nội dung.
\item Loại nội dung.
\item Người tạo nội dung.
\item Trạng thái nội dung.
\item Thời gian tạo và cập nhật.
\item Thông tin phục vụ tìm kiếm và tổ chức kho kiến thức cộng đồng.
\end{itemize}

Community Content có thể bao gồm educational articles, tutorials, equipment reviews, darkroom techniques, best practices và discussions.

\subsection{Workshop}

Hệ thống cần quản lý thông tin về các Workshop và training session được tổ chức trên nền tảng.

Các yêu cầu dữ liệu chính:

\begin{itemize}
\item Thông tin Workshop.
\item Chủ đề và nội dung.
\item Người tổ chức.
\item Thời gian tổ chức.
\item Địa điểm.
\item Số lượng người tham gia tối đa.
\item Chi phí tham gia nếu có.
\item Trạng thái Workshop.
\item Thông tin về người tham gia.
\end{itemize}

Dữ liệu Workshop phải hỗ trợ việc tổ chức, đăng ký và quản lý số lượng người tham gia.

\subsection{Photo}

Hệ thống cần quản lý thông tin về các ảnh kỹ thuật số thuộc bộ sưu tập của người dùng.

Các yêu cầu dữ liệu chính:

\begin{itemize}
\item Thông tin nhận diện và mô tả ảnh.
\item Tên hoặc tiêu đề ảnh.
\item Thông tin tham chiếu đến tệp ảnh.
\item Người quản lý ảnh.
\item Thời điểm thêm ảnh.
\item Thời điểm cập nhật ảnh.
\end{itemize}

Dữ liệu Photo phục vụ việc quản lý bộ sưu tập ảnh kỹ thuật số và có thể được sử dụng trong các chức năng image restoration và image enhancement của hệ thống.

\subsection{Complaint}

Hệ thống cần quản lý các khiếu nại và tranh chấp phát sinh trong quá trình sử dụng nền tảng.

Các yêu cầu dữ liệu chính:

\begin{itemize}
\item Người gửi khiếu nại.
\item Nội dung và chủ đề khiếu nại.
\item Đối tượng hoặc giao dịch liên quan.
\item Trạng thái xử lý.
\item Thời điểm tạo khiếu nại.
\item Thời điểm giải quyết.
\item Nội dung và kết quả xử lý.
\end{itemize}

Complaint phục vụ Administrator trong việc tiếp nhận, theo dõi và giải quyết tranh chấp giữa các bên.

\section{Quy tắc nghiệp vụ}

\subsection{Quy tắc về người dùng và vai trò}

Hệ thống quản lý người dùng thông qua thực thể \textbf{User}, mỗi tài khoản được gán một trong bốn vai trò sau:

\begin{itemize}
    \item \textbf{Photographer}: có thể tìm kiếm, đặt Creative Space, Resource hoặc Service Package; thực hiện thanh toán; tham gia Workshop; quản lý Photo cá nhân; viết Review và gửi Complaint.
    \item \textbf{Service Provider}: có thể tạo và quản lý Creative Space, Resource, Service Package và Promotion; xem lịch đặt chỗ và phiên sử dụng dịch vụ.
    \item \textbf{Photography Expert}: có thể tạo Community Content và tổ chức Workshop.
    \item \textbf{Administrator}: quản lý toàn bộ hệ thống, phê duyệt tài khoản Service Provider và xử lý Complaint.
\end{itemize}

Mỗi User chỉ có một vai trò tại một thời điểm. Tài khoản có thể ở trạng thái hoạt động hoặc bị khoá tuỳ theo quyết định của Administrator.

\subsection{Quy tắc về nhà cung cấp và tài nguyên}

\begin{itemize}
    \item Một \textbf{Service Provider} có thể quản lý nhiều \textbf{Creative Space}, nhiều \textbf{Resource}, nhiều \textbf{Service Package} và nhiều \textbf{Promotion}. Mỗi đối tượng trong số này chỉ thuộc về đúng một Service Provider.
    \item Mỗi \textbf{Resource} có thể trải qua nhiều lần bảo trì (\textbf{Maintenance}) trong vòng đời của nó. Mỗi Maintenance chỉ gắn với một Resource duy nhất.
    \item Khi một Resource đang trong trạng thái bảo trì, Resource đó không được phân bổ cho bất kỳ Reservation nào.
    \item Thông tin về số lượng (\texttt{quantity}), tình trạng (\texttt{condition}) và trạng thái (\texttt{status}) của Resource phải được cập nhật kịp thời để phản ánh đúng khả dụng thực tế.
\end{itemize}

\subsection{Quy tắc về dịch vụ}

\begin{itemize}
    \item Một \textbf{Service Package} có thể bao gồm nhiều \textbf{Creative Space} và nhiều \textbf{Resource}; ngược lại, một Creative Space hoặc Resource cũng có thể thuộc nhiều Service Package (quan hệ N:N).
    \item Khi tạo Service Package, Service Provider phải chỉ định ít nhất một Creative Space hoặc một Resource.
    \item Một \textbf{Promotion} do Service Provider tạo ra, có thời gian hiệu lực xác định (\texttt{start\_at}, \texttt{end\_at}) và chỉ được áp dụng trong khoảng thời gian đó.
    \item Promotion có thể thuộc loại giảm theo phần trăm hoặc giảm theo số tiền cố định, được xác định qua \texttt{discount\_type} và \texttt{discount\_value}.
\end{itemize}

\subsection{Quy tắc về đặt chỗ và sử dụng dịch vụ}

\begin{itemize}
    \item Một \textbf{Reservation} được tạo bởi đúng một User có vai trò Photographer. Một User có thể tạo nhiều Reservation.
    \item Một Reservation có thể đặt nhiều Creative Space và nhiều Resource đồng thời (quan hệ N:N). Reservation cũng có thể sử dụng một Service Package hoặc không.
    \item Trước khi xác nhận Reservation, hệ thống phải kiểm tra khả dụng của Creative Space và Resource trong khoảng thời gian \texttt{start\_time} đến \texttt{end\_time}. Nếu phát hiện xung đột lịch, Reservation không được xác nhận.
    \item Thời gian bắt đầu phải nhỏ hơn thời gian kết thúc: $\texttt{start\_time} < \texttt{end\_time}$.
    \item Một Reservation có thể có tối đa một \textbf{Payment}. Payment chỉ được tạo sau khi Reservation được xác nhận.
    \item Một Reservation có thể có tối đa một \textbf{Service Session}. Service Session chỉ được tạo sau khi Payment hoàn tất.
    \item Trong Service Session, \texttt{check\_out} phải lớn hơn \texttt{check\_in}; \texttt{actual\_usage\_duration} được tính từ hiệu số đó.
\end{itemize}

\subsection{Quy tắc về cộng đồng}

\begin{itemize}
    \item Một \textbf{User} có thể tạo nhiều \textbf{Community Content}. Mỗi Community Content chỉ thuộc về một User.
    \item Community Content có thể thuộc các loại: article, tutorial, technique, best practice hoặc discussion, được xác định qua thuộc tính \texttt{content\_type}.
    \item Một User có vai trò Photography Expert có thể tổ chức nhiều \textbf{Workshop}. Mỗi Workshop được tổ chức bởi đúng một User.
    \item Một User có thể tham gia nhiều Workshop; một Workshop có thể có nhiều User tham gia (quan hệ N:N). Số lượng người tham gia không được vượt quá \texttt{capacity} của Workshop.
    \item Một User, đặc biệt là Photographer, có thể quản lý nhiều \textbf{Photo}. Mỗi Photo chỉ thuộc về một User.
\end{itemize}

\subsection{Quy tắc về đánh giá và khiếu nại}

\begin{itemize}
    \item Một \textbf{Review} chỉ có thể được tạo sau khi Reservation tương ứng đã hoàn thành. Mỗi Reservation có tối đa một Review.
    \item Review phải được viết bởi chính User đã tạo Reservation đó.
    \item Một \textbf{Complaint} được tạo bởi một User. Một User có thể gửi nhiều Complaint; mỗi Complaint chỉ thuộc về một User.
    \item Complaint phải được Administrator tiếp nhận và xử lý. Khi xử lý xong, trường \texttt{resolved\_at} và \texttt{resolution} được cập nhật, trạng thái chuyển sang đã giải quyết.
\end{itemize}

\subsection{Các ràng buộc về thời gian và khả dụng}

\begin{itemize}
    \item \textbf{Ràng buộc thời gian Reservation}: \texttt{start\_time} $<$ \texttt{end\_time}; thời gian đặt không được nằm trong quá khứ tại thời điểm tạo Reservation.
    \item \textbf{Ràng buộc khả dụng Creative Space}: Tại cùng một khoảng thời gian, một Creative Space không thể xuất hiện trong hai Reservation đang hoạt động cùng lúc.
    \item \textbf{Ràng buộc khả dụng Resource}: Tổng số lượng Resource được phân bổ cho các Reservation đang hoạt động không được vượt quá \texttt{quantity} hiện có.
    \item \textbf{Ràng buộc Promotion}: Promotion chỉ có hiệu lực khi \texttt{start\_at} $\leq$ thời điểm đặt $\leq$ \texttt{end\_at} và trạng thái là đang hoạt động.
    \item \textbf{Ràng buộc Workshop}: Số lượng User đăng ký tham gia không được vượt quá \texttt{capacity}; \texttt{start\_time} $<$ \texttt{end\_time}.
    \item \textbf{Ràng buộc Maintenance}: Khoảng thời gian bảo trì của một Resource không được trùng với khoảng thời gian Resource đó được phân bổ trong Reservation.
\end{itemize}

\section{Tổng hợp yêu cầu và phạm vi dữ liệu}

\subsection{Danh sách đối tượng dữ liệu}

Hệ thống quản lý 15 thực thể chính và 6 bảng kết hợp như sau:

\begin{table}[H]
\centering
\caption{Danh sách thực thể của hệ thống}
\begin{tabular}{|c|l|p{9cm}|}
\hline
\textbf{STT} & \textbf{Thực thể} & \textbf{Ý nghĩa nghiệp vụ} \\
\hline
1  & User              & Đại diện người dùng của hệ thống với các vai trò Photographer, Service Provider, Photography Expert hoặc Administrator. \\
\hline
2  & Service Provider  & Đại diện cá nhân hoặc tổ chức cung cấp và quản lý creative space, resource và service package. \\
\hline
3  & Creative Space    & Đại diện các không gian sáng tạo như darkroom và photography studio. \\
\hline
4  & Resource          & Đại diện các tài nguyên như thiết bị chụp ảnh, thiết bị darkroom và vật tư tiêu hao. \\
\hline
5  & Maintenance       & Đại diện các hoạt động và lịch bảo trì tài nguyên. \\
\hline
6  & Service Package   & Đại diện các gói dịch vụ kết hợp creative space, resource và dịch vụ hỗ trợ. \\
\hline
7  & Promotion         & Đại diện các chương trình khuyến mãi do Service Provider cấu hình. \\
\hline
8  & Reservation       & Đại diện nghiệp vụ đặt creative space, resource hoặc service package. \\
\hline
9  & Payment           & Đại diện thông tin thanh toán phát sinh từ reservation. \\
\hline
10 & Service Session   & Đại diện phiên sử dụng dịch vụ thực tế bao gồm check-in, check-out. \\
\hline
11 & Review            & Đại diện đánh giá và nhận xét của người dùng sau quá trình sử dụng. \\
\hline
12 & Community Content & Đại diện các nội dung được chia sẻ trong cộng đồng nhiếp ảnh. \\
\hline
13 & Workshop          & Đại diện các buổi workshop hoặc training session. \\
\hline
14 & Photo             & Đại diện các ảnh kỹ thuật số được người dùng quản lý trên hệ thống. \\
\hline
15 & Complaint         & Đại diện các khiếu nại hoặc tranh chấp phát sinh trong quá trình sử dụng dịch vụ. \\
\hline
\end{tabular}
\end{table}

\begin{table}[H]
\centering
\caption{Danh sách bảng kết hợp}
\begin{tabular}{|c|l|l|p{6cm}|}
\hline
\textbf{STT} & \textbf{Mối kết hợp} & \textbf{Bảng kết hợp} & \textbf{Ý nghĩa nghiệp vụ} \\
\hline
1 & Service Package -- Creative Space & Package Space         & Liên kết các Creative Space được bao gồm trong Service Package. \\
\hline
2 & Service Package -- Resource       & Package Resource      & Lưu số lượng Resource được sử dụng trong một Service Package. \\
\hline
3 & Creative Space -- Reservation     & Reservation Space     & Liên kết các Creative Space được sử dụng trong Reservation. \\
\hline
4 & Reservation -- Resource           & Reservation Resource  & Lưu số lượng Resource được phân bổ cho một Reservation. \\
\hline
5 & Service Session -- Resource       & Session Resource      & Lưu số lượng Resource thực tế được sử dụng trong một Service Session. \\
\hline
6 & User -- Workshop                  & Workshop Registration & Liên kết User với các Workshop mà User tham gia. \\
\hline
\end{tabular}
\end{table}

\subsection{Danh sách yêu cầu dữ liệu}

Các yêu cầu dữ liệu được xác định dựa trên các nhóm nghiệp vụ chính của hệ thống:

\begin{itemize}
    \item \textbf{Quản lý người dùng}: lưu thông tin tài khoản, vai trò, trạng thái và thông tin cơ bản của người dùng.
    \item \textbf{Quản lý nhà cung cấp và tài nguyên}: quản lý Service Provider, Creative Space, Resource, Maintenance, Service Package và Promotion.
    \item \textbf{Quản lý đặt và sử dụng dịch vụ}: quản lý Reservation, Payment và Service Session nhằm theo dõi toàn bộ vòng đời của một lần sử dụng dịch vụ.
    \item \textbf{Quản lý cộng đồng}: quản lý Community Content, Workshop, Photo và các hoạt động chia sẻ kiến thức.
    \item \textbf{Quản lý phản hồi và khiếu nại}: quản lý Review và Complaint để hỗ trợ đánh giá chất lượng dịch vụ và xử lý tranh chấp.
\end{itemize}

Các dữ liệu cần được lưu trữ lâu dài bao gồm:
\begin{itemize}
    \item \textbf{Dữ liệu chủ}: User, Service Provider, Creative Space, Resource, Service Package.
    \item \textbf{Dữ liệu giao dịch}: Reservation, Payment, Service Session, Maintenance, Review, Promotion, Complaint.
    \item \textbf{Dữ liệu nội dung}: Community Content, Workshop, Photo.
\end{itemize}

\subsection{Danh sách quy tắc nghiệp vụ}

\begin{table}[H]
\centering
\caption{Tổng hợp các mối quan hệ và bản số}
\begin{tabular}{|c|l|l|c|}
\hline
\textbf{STT} & \textbf{Mối kết hợp} & \textbf{Quan hệ} & \textbf{Bản số} \\
\hline
1  & Service Provider -- Creative Space  & Quản lý          & 1:N    \\
\hline
2  & Service Provider -- Resource        & Quản lý          & 1:N    \\
\hline
3  & Service Provider -- Service Package & Tạo              & 1:N    \\
\hline
4  & Service Provider -- Promotion       & Tạo              & 1:N    \\
\hline
5  & Resource -- Maintenance             & Được bảo trì     & 1:N    \\
\hline
6  & Service Package -- Creative Space   & Bao gồm          & N:N    \\
\hline
7  & Service Package -- Resource         & Bao gồm          & N:N    \\
\hline
8  & User -- Reservation                 & Đặt              & 1:N    \\
\hline
9  & Creative Space -- Reservation       & Được đặt         & N:N    \\
\hline
10 & Service Package -- Reservation      & Được chọn        & N:1    \\
\hline
11 & Reservation -- Resource             & Phân bổ          & N:N    \\
\hline
12 & Reservation -- Payment              & Có thanh toán    & 1:0..1 \\
\hline
13 & Reservation -- Service Session      & Có phiên sử dụng & 1:0..1 \\
\hline
14 & Service Session -- Resource         & Sử dụng thực tế  & N:N    \\
\hline
15 & User -- Review                      & Viết             & 1:N    \\
\hline
16 & Reservation -- Review               & Được đánh giá    & 1:0..1 \\
\hline
17 & User -- Community Content           & Tạo              & 1:N    \\
\hline
18 & User -- Workshop                    & Tổ chức          & 1:N    \\
\hline
19 & User -- Workshop                    & Tham gia         & N:N    \\
\hline
20 & User -- Photo                       & Quản lý          & 1:N    \\
\hline
21 & User -- Complaint                   & Gửi              & 1:N    \\
\hline
\end{tabular}
\end{table}

\subsection{Đầu vào cho thiết kế quan niệm}

Dựa trên kết quả phân tích yêu cầu, các thành phần đầu vào cho giai đoạn thiết kế quan niệm bao gồm:

\begin{itemize}
    \item \textbf{Danh sách 15 thực thể} đã được xác định cùng với các thuộc tính và khoá chính tương ứng.
    \item \textbf{Danh sách 21 mối quan hệ} giữa các thực thể, bao gồm tên quan hệ, bản số và ràng buộc tham gia.
    \item \textbf{Danh sách 6 bảng kết hợp} cho các quan hệ N:N, một số có thuộc tính riêng như \texttt{quantity}.
    \item \textbf{Các quy tắc nghiệp vụ và ràng buộc} về thời gian, khả dụng, vai trò và trạng thái.
    \item \textbf{Từ điển dữ liệu} mô tả ý nghĩa nghiệp vụ của từng thuộc tính trong mỗi thực thể.
\end{itemize}

Toàn bộ thông tin trên sẽ được sử dụng làm nền tảng để xây dựng sơ đồ EER và thiết kế cơ sở dữ liệu ở mức quan niệm.